% !TEX program = xelatex
\documentclass[12pt,a4paper,UTF8]{ctexart}

\usepackage{amsmath,amssymb,mathtools,amsthm}
\usepackage{geometry}
\geometry{left=2.6cm,right=2.6cm,top=2.6cm,bottom=2.6cm}
\usepackage{enumitem}
\usepackage{booktabs}
\usepackage{hyperref}
\hypersetup{colorlinks=true,linkcolor=blue,urlcolor=blue,citecolor=blue}

\newcommand{\Z}{\mathbb{Z}}
\newcommand{\R}{\mathbb{R}}
\newcommand{\C}{\mathbb{C}}
\newcommand{\Lie}{\mathfrak{g}}
\newcommand{\gl}{\mathfrak{gl}}
\newcommand{\su}{\mathfrak{su}}
\newcommand{\so}{\mathfrak{so}}
\newcommand{\ad}{\operatorname{ad}}
\DeclareMathOperator{\Ad}{Ad}
\DeclareMathOperator{\End}{End}
\DeclareMathOperator{\Aut}{Aut}
\DeclareMathOperator{\GL}{GL}
\DeclareMathOperator{\SL}{SL}
\DeclareMathOperator{\SO}{SO}
\DeclareMathOperator{\SU}{SU}
\DeclareMathOperator{\Sp}{Sp}
\DeclareMathOperator{\tr}{tr}
\DeclareMathOperator{\rank}{rank}
\DeclareMathOperator{\image}{im}

\newtheorem{theorem}{定理}[section]
\newtheorem{lemma}[theorem]{引理}
\newtheorem{proposition}[theorem]{命题}
\newtheorem{corollary}[theorem]{推论}
\theoremstyle{definition}
\newtheorem{definition}[theorem]{定义}
\newtheorem{example}[theorem]{例}
\newtheorem{exercise}[theorem]{习题}
\newtheorem{remark}[theorem]{注}

\title{代数拓扑极速入门：拓扑空间、连续映射与范畴}
\author{HUST SFT 周炜皓}
\date{}

\begin{document}
\maketitle

\begin{abstract}
本讲义试图带你极速入门代数拓扑，并尽可能只介绍必要的前置知识，同时穿插阐述代数拓扑的核心思想。只要求读者有集合论基础，如果知道一点群论和点集拓扑会更舒适。内容包括：拓扑空间与连续映射、同胚、子空间/积空间/商空间三种基本构造、连通性与紧致性，以及范畴与函子这一组织代数拓扑的整体语言。
\end{abstract}

\tableofcontents
\newpage

%============================================================
\section{从几何到拓扑\\
From geometry to topology}
%============================================================

代数拓扑研究如何利用代数工具刻画拓扑空间的结构。其基本思想是：将几何空间转化为群、环、模等代数对象，再通过代数性质研究空间本身。

例如，圆周 $S^1$、球面 $S^2$ 和环面 $T^2$ 是三种不同的拓扑空间。我们希望回答：如何严格区分它们？如何描述空间中的洞？如何判断一个空间能否连续变形为另一个空间？

为此，我们首先需要明确什么是拓扑空间，以及什么是空间之间的连续映射。

%============================================================
\section{拓扑空间\\
Topological spaces}
%============================================================

\begin{definition}[拓扑 Topology]
设 $X$ 为集合。若集合族 $\mathcal T\subseteq\mathcal P(X)$ 满足：
\begin{enumerate}[label=(\arabic*),itemsep=2pt]
\item $\varnothing,X\in\mathcal T$；
\item 任意多个 $\mathcal T$ 中集合的并仍属于 $\mathcal T$；
\item 有限多个 $\mathcal T$ 中集合的交仍属于 $\mathcal T$，
\end{enumerate}
则称 $\mathcal T$ 为 $X$ 上的一个拓扑，二元组 $(X,\mathcal T)$ 称为拓扑空间。集合族 $\mathcal T$ 中的元素称为开集。
\end{definition}

拓扑公理刻画了一种抽象的邻近关系。它不要求空间具有距离，只要求开集满足上述三个条件。特别地，并集允许任意多个，而交集只要求有限多个。

\begin{remark}
为什么交集只能取有限多个？在 $\R$ 的通常拓扑下，
\[
\bigcap_{n=1}^{\infty}\Bigl(-\frac1n,\frac1n\Bigr)=\{0\}
\]
并不是开集。若允许任意交，则单点集将成为开集，离散化会使拓扑失去区分能力。有限交封闭恰好对应“有限多个条件同时成立仍应是开条件”。
\end{remark}

\begin{example}[典型拓扑]
任意集合 $X$ 上存在两种极端的拓扑：离散拓扑 $\mathcal T=\mathcal P(X)$，其中所有子集都是开集；平凡拓扑 $\mathcal T=\{\varnothing,X\}$，其中只有空集与全集是开集。实数集 $\R$ 上通常的拓扑则由所有开区间生成，其开集均可表示为开区间的任意并。
\end{example}

\begin{definition}[闭集 Closed set]
若 $F\subseteq X$ 的补集 $X\setminus F$ 是开集，则称 $F$ 为闭集。由拓扑公理可知，任意多个闭集的交以及有限多个闭集的并仍为闭集。注意，开集与闭集并不互斥，一个集合可以同时是开集和闭集，称为既开又闭集（clopen set）。
\end{definition}

%============================================================
\section{连续映射与同胚\\
Continuous maps and homeomorphisms}
%============================================================

拓扑学研究空间之间的连续变形，因此需要一种不依赖距离的连续性定义。

\begin{definition}[连续映射 Continuous map]
设 $X,Y$ 为拓扑空间，$f:X\to Y$ 为映射。若对任意开集 $V\subseteq Y$，其逆像 $f^{-1}(V)$ 在 $X$ 中开，则称 $f$ 连续。
\end{definition}

这里需要特别注意：连续性要求开集的逆像为开集，而不要求开集的像为开集。连续映射具有以下基本性质：恒等映射连续，连续映射的复合仍连续，连续映射在子空间上的限制仍连续。

\begin{definition}[同胚 Homeomorphism]
若映射 $f:X\to Y$ 是双射，且 $f$ 与 $f^{-1}$ 均连续，则称 $f$ 为同胚，并称 $X,Y$ 同胚，记作 $X\cong Y$。
\end{definition}

同胚是拓扑空间之间最基本的等价关系。两个空间同胚，意味着它们具有完全相同的拓扑结构，只是具体表示不同。例如，圆周与正方形边界同胚，开区间 $(0,1)$ 与实数轴 $\R$ 同胚。

因此，拓扑学关注的不是空间的具体形状，而是其在同胚意义下保持不变的性质。

%============================================================
\section{常见的空间构造\\
Subspaces, products and quotients}
%============================================================

代数拓扑中经常需要从已有空间构造新的空间。最重要的三种构造是子空间、积空间和商空间。

\begin{definition}[子空间 Subspace]
设 $A\subseteq X$，其中 $X$ 为拓扑空间。定义
\[
\mathcal T_A=\{A\cap U:U\in\mathcal T_X\},
\]
则 $\mathcal T_A$ 是 $A$ 上的拓扑，称为子空间拓扑。由此得到的拓扑空间称为 $X$ 的子空间。
\end{definition}

例如，$[0,1)$ 在 $[0,1]$ 的子空间拓扑下是开集，尽管它不是实数轴 $\R$ 中的开集。这说明开集的概念始终依赖于所在的拓扑空间。

\begin{definition}[积拓扑 Product topology]
设 $X,Y$ 为拓扑空间。$X\times Y$ 上由所有开矩形 $U\times V$（其中 $U\subseteq X,V\subseteq Y$ 分别为开集）生成的拓扑称为积拓扑。
\end{definition}

例如，$\R^2=\R\times\R$ 的积拓扑就是通常的欧氏拓扑。更一般地，任意拓扑空间族都可以构造其积空间，积拓扑由有限多个坐标方向上的开集限制生成。

\begin{definition}[商拓扑 Quotient topology]
设 $X$ 为拓扑空间，$\sim$ 为其上的等价关系，记商集为 $X/{\sim}$，自然投影为
\[
q:X\to X/{\sim},\qquad q(x)=[x].
\]
定义 $U\subseteq X/{\sim}$ 为开集，当且仅当 $q^{-1}(U)$ 在 $X$ 中开。由此得到的拓扑称为商拓扑，所得空间称为商空间。商拓扑的几何意义是：将原空间中的某些点识别为同一个点，同时尽可能保留原有的拓扑结构。
\end{definition}

\begin{example}[圆周的构造]
将闭区间 $[0,1]$ 的两个端点识别，即规定 $0\sim1$，其余点各自构成等价类，则所得商空间同胚于圆周 $S^1$。事实上，映射
\[
f:[0,1]\to S^1,\qquad f(t)=e^{2\pi i t},
\]
诱导出这一同胚。
\end{example}

\begin{proposition}[商空间的泛性质 Universal property]
设 $q:X\to Y$ 为商映射，$Z$ 为拓扑空间，则映射 $f:Y\to Z$ 连续，当且仅当 $f\circ q:X\to Z$ 连续。
\end{proposition}

\begin{proof}
若 $f$ 连续，则 $f\circ q$ 作为连续映射的复合连续。反之设 $f\circ q$ 连续，取 $Z$ 中开集 $W$，则 $(f\circ q)^{-1}(W)=q^{-1}\bigl(f^{-1}(W)\bigr)$ 在 $X$ 中开；由商拓扑的定义，$f^{-1}(W)$ 在 $Y$ 中开，故 $f$ 连续。
\end{proof}

这个性质使得我们可以通过原空间上的连续性来研究商空间上的映射。

%============================================================
\section{连通性与紧致性\\
Connectedness and compactness}
%============================================================

连通性和紧致性是拓扑空间最重要的基本性质，也是后续研究流形、同伦与同调时经常使用的工具。

\begin{definition}[连通空间 Connected space]
若拓扑空间 $X$ 不能表示为两个非空、不相交的开集之并，则称 $X$ 连通。等价地，$X$ 不存在非平凡的既开又闭子集。
\end{definition}

\begin{definition}[道路连通空间 Path-connected space]
若对任意 $x,y\in X$，存在连续映射 $\gamma:[0,1]\to X$，使得 $\gamma(0)=x,\ \gamma(1)=y$，则称 $X$ 道路连通。
\end{definition}

道路连通必然蕴含连通，但反之不成立。道路连通描述了空间中任意两点能否由连续道路连接，而连通性只要求空间不能被分割成两个不相交的非空开集。经典的反例是“拓扑学家的正弦曲线”
\[
S=\Bigl\{(x,\sin\tfrac1x):x>0\Bigr\}\cup\bigl(\{0\}\times[-1,1]\bigr),
\]
它连通但不道路连通。

\begin{definition}[紧致空间 Compact space]
若拓扑空间 $X$ 的任意开覆盖都存在有限子覆盖，则称 $X$ 紧致。这里开覆盖是指一族开集 $\{U_\lambda\}_{\lambda\in\Lambda}$，满足 $X=\bigcup_{\lambda\in\Lambda}U_\lambda$。同理可以定义紧集。
\end{definition}

在欧氏空间中，Heine--Borel 定理表明：$\R^n$ 的子集紧致，当且仅当它闭且有界。但这一刻画依赖于欧氏空间结构，在一般拓扑空间中并不成立。

%============================================================
\section{范畴与函子：代数拓扑的语言\\
Categories and functors}
%============================================================

前面的内容主要研究拓扑空间本身。现在引入范畴论，它将成为我们理解代数拓扑整体结构的重要工具。

\begin{definition}[范畴 Category]
一个范畴 $\mathcal C$ 由以下数据组成：一族对象 $\operatorname{Ob}(\mathcal C)$；对每对对象 $X,Y$，一个态射集 $\operatorname{Hom}_{\mathcal C}(X,Y)$；以及态射的复合运算
\[
\operatorname{Hom}_{\mathcal C}(Y,Z)\times\operatorname{Hom}_{\mathcal C}(X,Y)\to\operatorname{Hom}_{\mathcal C}(X,Z),\qquad (g,f)\mapsto g\circ f,
\]
并满足复合的结合律，且对每个对象 $X$ 存在恒等态射 $\operatorname{id}_X\in\operatorname{Hom}_{\mathcal C}(X,X)$ 满足单位律。对象表示所研究的数学结构，态射表示结构之间保持相应性质的映射。
\end{definition}

例如，集合与集合映射构成范畴 $\mathbf{Set}$；群与群同态构成范畴 $\mathbf{Grp}$；拓扑空间与连续映射构成范畴 $\mathbf{Top}$；Abel 群与群同态构成范畴 $\mathbf{Ab}$。

在 $\mathbf{Top}$ 中，对象是拓扑空间，态射是连续映射。两个空间同胚，当且仅当它们在 $\mathbf{Top}$ 中同构。

\begin{definition}[函子 Functor]
设 $\mathcal C,\mathcal D$ 为范畴。协变函子 $F:\mathcal C\to\mathcal D$ 将 $\mathcal C$ 中的对象 $X$ 对应到 $\mathcal D$ 中的对象 $F(X)$，并将态射 $f:X\to Y$ 对应到态射 $F(f):F(X)\to F(Y)$，且满足
\[
F(\operatorname{id}_X)=\operatorname{id}_{F(X)},\qquad F(g\circ f)=F(g)\circ F(f).
\]
若态射方向反转，即 $F(f):F(Y)\to F(X)$，则称为反变函子。
\end{definition}

这个定义和群同态其实很像。函子的意义在于，它能够将一种数学结构系统地转化为另一种数学结构，同时保留态射的复合关系。

这正是代数拓扑的核心思想。给定拓扑空间 $X$，我们可以构造基本群 $\pi_1(X)$、同调群 $H_n(X)$ 和上同调群 $H^n(X)$。更重要的是，连续映射也会诱导相应的代数映射。

\begin{remark}[协变与反变]
基本群和同调具有协变性：连续映射 $f:X\to Y$ 诱导
\[
\pi_1(f):\pi_1(X)\to\pi_1(Y),\qquad H_n(f):H_n(X)\to H_n(Y)
\]
（基本群的构造需要指定基点）。而上同调具有反变性：$f$ 诱导
\[
f^*:H^n(Y)\to H^n(X).
\]
因此，代数拓扑并非简单地为每个空间计算一个代数不变量，而是研究拓扑空间及其映射如何系统地转化为代数结构及其映射。函子性保证了“空间同胚 $\Rightarrow$ 代数对象同构”，这正是用代数不变量区分空间的合法性来源。
\end{remark}

%============================================================
\section{习题\\
Exercises}
%============================================================

\begin{exercise}
证明离散拓扑和平凡拓扑均满足拓扑公理。
\end{exercise}

\begin{exercise}
证明同胚关系是拓扑空间之间的等价关系。
\end{exercise}

\begin{exercise}
证明连通空间的连续像仍然连通。
\end{exercise}

\begin{exercise}
证明道路连通空间必然连通，并举例说明逆命题不成立。
\end{exercise}

\begin{exercise}
证明商映射 $q:X\to Y$ 满足：$f:Y\to Z$ 连续，当且仅当 $f\circ q$ 连续。
\end{exercise}

%============================================================
\section{部分习题提示与解答\\
Hints and solutions}
%============================================================

\begin{enumerate}[label=\textbf{\arabic*.},itemsep=4pt]
\item 略。
\item 只需验证等价关系的三条：自反性取恒等映射；对称性取 $f^{-1}$；传递性取同胚的复合。
\item 反证法：设 $f(X)=A\cup B$ 为两个非空不交开集之并，则 $f^{-1}(A),f^{-1}(B)$ 是 $X$ 中非空不交开集，与 $X$ 连通矛盾。
\item $[0,1]$ 是连通的，利用道路连通的定义和上一题的结论即可：任取 $x,y$，存在道路 $\gamma:[0,1]\to X$ 连接它们，故 $X=\bigcup_{\gamma}\gamma([0,1])$ 为一族含基点的连通集之并，从而是连通的。反例见正文中拓扑学家的正弦曲线。
\item 利用连续的定义和商拓扑的定义即可：一边由复合的连续性直接得到，另一边用 $q^{-1}(f^{-1}(W))$ 为开集推出 $f^{-1}(W)$ 为开集。
\end{enumerate}

%============================================================
\section*{参考文献}
\addcontentsline{toc}{section}{参考文献}
%============================================================

\begin{thebibliography}{9}
\bibitem{hatcher} Allen Hatcher, \emph{Algebraic Topology}, Chapter 0.
\bibitem{munkres-top} James R. Munkres, \emph{Topology}, Chapters 2--4.
\bibitem{munkres-alg} James R. Munkres, \emph{Elements of Algebraic Topology}.
\bibitem{may} J. Peter May, \emph{A Concise Course in Algebraic Topology}.
\end{thebibliography}

\end{document}